% VH-exists paper: a dated record of what is solid. SKELETON (outline stage), 6 Oct 2026.
% Every claim carries one label: [hand] [compiled] [computed] [cited] [open].
% No claim stronger than the Navigator's ledger. Author: Kyle Mathewson (the user's statement, relayed 12:06). Acknowledgement of AI assistance: full disclosure section (user's decision 12:20, 6 Oct 2026); Kyle approves the final wording before posting.
\documentclass[11pt]{article}
\usepackage{amsmath,amssymb,amsthm,hyperref,url}
\newtheorem{theorem}{Theorem}\newtheorem{lemma}{Lemma}\newtheorem{conjecture}{Conjecture}
\newcommand{\lab}[1]{\textsf{[#1]}}
\title{The vacancy induction for the Four Colour Theorem:\\ a dated record of what is solid}
\author{Kyle Mathewson}
\date{6 October 2026 (draft outline)}
\begin{document}\maketitle
\begin{abstract}TODO after sections 1--6 are drafted. States only what the ledger supports.\end{abstract}
\section*{How this paper was made}
\textbf{The mathematics, the Lean proofs, the computations and the text of this paper were produced by AI agents.} No human mathematician wrote any of the proofs or the code. The agents were instances of Claude, a family of AI models made by Anthropic, running in Claude Code: separate, long-running sessions that could read and write files, run programs on the author's computers, and exchange dated messages through a shared repository. Each session had a fixed role.
\begin{itemize}
\item \emph{Long Table} proposed ideas, wrote the hand proofs and counterexamples, and designed and ran the declared experiments.
\item \emph{Math} reviewed hand proofs, accepted or refused results, wrote the Lean~4 formalisation and ran its module audit.
\item \emph{Audit} reran results with its own code, which imported nothing from the other teams, and read every claim adversarially.
\item \emph{Navigator} kept the status ledger and was the only session allowed to call a result proved, compiled, exploring or killed.
\item A \emph{coordinator} session sequenced work between the teams.
\item A \emph{night swarm} of overnight exploratory agents produced leads, which this paper uses only as leads.
\end{itemize}
Sessions often delegated sub-tasks to short-lived sub-agents, for example to write a checker for an experiment without seeing the program it checks. Several Claude model versions were used over the course of the project, and the version behind each piece of work was not recorded consistently [\textsc{Kyle: name versions or leave as is}].

\textbf{What the author did.} Kyle Mathewson directed the project and is responsible for every claim in it. From the message record, he chose the problem and the direction, set the standing rules (what experiments may run, under what limits, and who may release them), assigned the roles and the coordinator's authority, approved actions with outside effects such as installing software, redirected the work when it drifted (for example, from brute-force search back to new ideas), and decided authorship and this disclosure [\textsc{Kyle: correct or add in your own words}]. He did not write the proofs or the code.

\textbf{How the results were checked.} Every claim carries one label, and the label says how far it has been checked.
\begin{itemize}
\item \lab{compiled}: proved in Lean~4 with no \texttt{sorry} and only Lean's standard axioms, and included in the module audit. The Lean kernel, not an AI, is what certifies these.
\item \lab{hand}: a written proof reviewed by a different AI session. \emph{No human referee has checked any of these.} AI reviewers can share blind spots, so treat them as unrefereed.
\item \lab{computed}: an experiment whose statement, programs and limits were fixed and hashed before it ran, whose output was checked by a separately written program, and which claims only what happened on the graphs it ran.
\item \lab{cited}: taken from the published literature as summarised by the agents; consult the original.
\item \lab{open}: not claimed.
\end{itemize}
No status in this paper is stronger than the Navigator's ledger. Mistakes were made along the way---wrong timestamps, a mislabelled process, conjectures that were later refuted---and each was corrected by a dated message rather than erased.

\textbf{The record.} The repository at \url{https://github.com/kylemath/GraphColourDeductiveSolver} [\textsc{confirm it is public before posting}] holds every message between the teams, every declaration with its hash, and all code and outputs. Each statement in this paper can be traced to a dated file there.

\section{The vacancy induction and the hypothesis VH$\exists$}
TODO: definitions; theorem VH$\exists\Rightarrow$ 4-colourability; containment and apex-singleton lemmas. \lab{hand}/\lab{compiled}
\section{What is compiled in Lean}
TODO: short fills; three-move obstruction; mobility (general, triangulated); clique and protected lifts; unequal-pole belt walk; Lemma L4; Theorem P; belt vacancy theorem (only if audited); module audit counts; axioms. \lab{compiled}
\section{What a least failure must look like}
TODO: fixed-hole theorem; face-avoiding reduction; degree-6 split; Theorem H (hand, reviewed by Math, not audited). \lab{hand}
\section{Negative results}
TODO: Conjecture L refuted (W6; $A_r$ orbits that still fill in 2--3 swaps); radius $\le 3$ refuted (T4); killed fitted statements as lessons; pathway kills (P-B target, P-C). \lab{computed}/\lab{hand}
\section{Finite evidence, with exact scope}
TODO: WP18, WP19; WP20 P1 and WP21 only once their checks are complete, else ``pending''. Results ``on these graphs'' only. \lab{computed}
\section{Open problems}
TODO: VH$\exists$; Conjecture R; the pathways. \lab{open}
\end{document}
